\documentclass[12pt]{article}

% fonts

\usepackage[T1]{fontenc}
\usepackage{newtxtext}
\usepackage{bm}
\usepackage{amsthm} % before newtxmath b/c latex
\usepackage{newtxmath}
% \usepackage[bb=mt,cal=mt]{mathalpha}

\usepackage{amsmath}
\usepackage{amsfonts}
\usepackage{thmtools}
\usepackage{mathtools}

% spacing

\usepackage[parfill]{parskip}
\setlength{\parindent}{0pt}
\setlength{\parskip}{2ex plus 0.2ex minus 0.2ex}

% geometry of the page

\usepackage[
paper  = letterpaper,
left   = 1in,
right  = 1in,
top    = 1in,
bottom = 1in,
]{geometry}

% references

\usepackage{natbib}


\begin{document}

\begin{flushleft}
\textbf{A Stochastic Approximation Method (Robbins et al., 1951)} \\
Satvik Saha (ss7228) \\
\today
\end{flushleft}

\vspace{0.1in}


\cite{robbins51} describes a root finding algorithm which replaces direct
computation of the function with a (random) estimate.
Given a suitable increasing function $M$, one can solve for its root $x^*$
(satisfying $M(x^*) = 0$) by iterating $x_{n + 1} = x_n - a_n M(x_n)$ with an
appropriate step size sequence $\{a_n\}$.
% The root $x^*$ is indeed a fixed point of this procedure.
Assuming that $M$ is differentiable with $0 < \eta < M' < \gamma$ gives $M(x_n)
- M(x^*) = M'(x_n')(x_n - x^*)$ for some $x_n'$ between $x_n, x^*$, hence with
all $0 < a_n \leq 1 /\gamma$,
\begin{align*}
    |x_{n + 1} - x^*|
        &= |x_n - x^* - a_n M'(x_n')(x_n - x^*)| \\
        &= (1 - a_n M'(x_n')) |x_n - x^*| \\
        &\leq (1 - \eta a_n) |x_n - x^*| \\
        &\leq e^{-\eta a_n} |x_n - x^*| \\
        &\leq e^{-\eta\sum_{k = 1}^n a_n} |x_1 - x^*|.
\end{align*}
If $\{a_n\}$ is chosen such that $\sum_{n = 1}^\infty a_n = \infty$ (optimally,
$a_n = 1 /\gamma$), we have convergence with $x_n \to x^*$.

If computing $M(x)$ is intractable, we may instead use an estimate $y(x)$,
specifically an unbiased estimator satisfying $\mathbb{E}[y(x) \mid x] = M(x)$
and $\mathbb{E}[y(x)^2 \mid x]$ uniformly bounded.
Our new iterative scheme is $x_{n + 1} = x_n - a_n y_n$, where $y_n$ is an
instance of $y(x_n)$.
With some careful conditions on $M$ and $y$, \cite{robbins51} show that for
sequences $\{a_n\}$ of \emph{type} $1 /n$ (we want $c /n < a_n < c' /n$ for
constants $c, c' > 0$), we have convergence with $x_n \to x^*$ in $L^2$ (hence
in probability).

In the context of minimizing a function $f$, the gradient $\nabla{f}$ plays the
role of $M$ here; the algorithm is one of \emph{stochastic gradient descent}
(SGD).
The conditions on $M$ thus correspond to assumptions on the smoothness and
convexity of $f$.

We have seen how SGD can be used effectively to compute quantities like
\emph{maximum a posteriori} (MAP) estimates of parameters in probabilistic
models.
Here, the gradient of the objective function to be minimized (the negative log
joint density) is often expensive to compute since it depends on all of the
given data $x$.
Since an unbiased estimate of $\frac{1}{n}\sum_{i = 1}^n g(x_i)$ is $g(x_I)$
(or indeed $\frac{1}{m}\sum_{j = 1}^m g(x_{I_j})$), where $I$ (or each $I_j$)
is uniformly random over the indices $\{1, \dots, n\}$, SGD lets us deal with
gradients of this form using a small fraction of our data at each iterate.

Stochastic gradients are a key idea in \cite{ranganath14}: the gradient of the
ELBO to be optimized is an expectation over the variational distribution $q$,
which is estimated by a Monte-Carlo average using samples from $q$.


\nocite{*}
\bibliographystyle{plainnat}

\vfill

\begingroup
\renewcommand{\section}[2]{}%
\textbf{References}
\bibliography{week04}
\endgroup


\end{document}

%%% Local Variables:
%%% mode: latex
%%% TeX-master: t
%%% End:
